\chapter{Recherche des problèmes du code principal}

Ce chapitre se concentrera sur le travail réalisé durant ce travail de Bachelor.

\section{Vérifications des interruptions}

Afin de vérifier si l'interruption liée à la communication sans fil viendrait à gêner l'interruption principale (ADC, régulation, PWM,...), il suffit simplement de mettre en commentaire les parties concernant la communication sans fil et flashé le DSP afin de vérifier l'hypothèses.\\

Les parties mises en commentaire sont les suivantes :

\begin{minted}{c}
    __interrupt void INT_mySCI0_RX_ISR(void) // .xxx
    {
        uint16_t c;

        while (SCI_getRxFIFOStatus(SCIA_BASE) > 0) {
            c = SCI_readCharBlockingFIFO(SCIA_BASE);
            if (rxIndex < RX_BUF_LEN) {
                rxBuffer[rxIndex++] = c;
            }
            if(c == '\x02'){
                // Marque le début de trame
                dataIndex = rxIndex; // Sauvegarde la position pour les données
            }
            if(c == '\x03'){
                // Fin de trame - traite les données
                if(dataIndex > 0 && dataIndex < rxIndex-1) {
                    state_PIN = (rxBuffer[dataIndex] == '1') ? 1 : 0;
                }
                // Reset du buffer
                rxIndex = 0;
                int b;
                for(b = 0; b < RX_BUF_LEN; b++) {
                    rxBuffer[b] = 0;
                }
            }
        }

        SCI_clearOverflowStatus(SCIA_BASE);
        SCI_clearInterruptStatus(SCIA_BASE, SCI_INT_RXFF);
        Interrupt_clearACKGroup(INTERRUPT_ACK_GROUP9);
    }

    __interrupt void INT_mySCI0_TX_ISR(void)
    {
        while (SCI_getTxFIFOStatus(SCIA_BASE) < SCI_FIFO_TX16 && txIndex < txLength) {
            SCI_writeCharBlockingFIFO(SCIA_BASE, txBuffer[txIndex++]);
        }

        if (txIndex >= txLength) {
            SCI_disableInterrupt(mySCI0_BASE, SCI_INT_TXFF); // Fin d’envoi
        }

        // Acquitter l'interruption
        SCI_clearOverflowStatus(SCIA_BASE);
        SCI_clearInterruptStatus(SCIA_BASE, SCI_INT_TXFF);
        Interrupt_clearACKGroup(INTERRUPT_ACK_GROUP9); // For the PIE acknowledge
    }
\end{minted}

\section{Vérifications des temps de l'intéruption principal}

\section{Réecriture de la machine d'état pour la régulation}